\section{Principal hyperbolic quotients over every imaginary quadratic field}
\label{sec:principal-quotients}
This section reduces all split integral polarization types in a fixed dimension to one principal hyperbolic family, with a degree cost linear in the polarization volume.

Let $K$ be an imaginary quadratic field, let $\mathcal O=\mathcal O_K$, and define $t_K\in\{1,2\}$ by
\[
 \operatorname{Tr}_{K/\Q}(\mathcal O)=t_K\Z.
\]
Thus $t_K=1$ when the discriminant is odd and $t_K=2$ when it is even. For a projective $\mathcal O$-module $Q$, write
\[
 Q^\vee=\operatorname{Hom}_{\bar{\mathcal O}}(Q,\mathcal O).
\]
Its elements are additive functions satisfying $p(aq)=\bar a p(q)$. The hyperbolic Hermitian module $H(Q)$ is $Q^\vee\oplus Q$, with pairing
\[
 ((p,q),(p',q'))\longmapsto p(q')+\overline{p'(q)}.
\]
We use Hermitian forms linear in the first variable.

\begin{theorem}\label{thm:principal-quotient}
Let $(A,L)$ be a split Weil $2n$-fold over $K$, with compatible polarization and integral action by $\mathcal O_f=\Z+f\mathcal O$. Put $\Delta=\chi(L)$. There is a $K$-linear isogeny to a principally polarized split Weil variety,
\[
 q:A\longrightarrow(P,\Theta),\qquad
 q^*\Theta=t_K fL,\qquad
 \deg q=(t_K f)^{2n}\Delta,
\]
where equality of polarizations means equality of their first Chern classes. The integral Hermitian lattice of $P$ is $H(Q)$ for a projective rank-$n$ module $Q$.

At most $h_K$ isomorphism types of such hyperbolic modules occur, where $h_K$ is the ideal class number. If $C_I$ bounds detectors throughout the principal family associated with $H(\mathcal O^{n-1}\oplus I)$, then
\[
 \delta_W(A,L)\le t_K^n f^n\Delta\,
       \max_{[I]\in\operatorname{Cl}(\mathcal O)} C_I.
\]
The same conclusion on CM fibres requires the bounds only on the CM points of those principal families. If $\mathcal O$ has class number one, only the lattice $U^n$ is needed.

More generally, put
\[
 m_K=\max_{[I]\in\operatorname{Cl}(\mathcal O)}
       \min_{\substack{J\subset\mathcal O\\{[J]=[I]}}}
                    \operatorname{Nm}(J).
\]
If $C_U$ bounds detectors on the single principal family with lattice $U^n$, then
\[
 \delta_W(A,L)\le(t_Kfm_K)^n\Delta C_U,
 \qquad m_K\le\frac2\pi\sqrt{|\operatorname{disc}K|}.
\]
Again the CM version holds with a supplier only on CM points.
\end{theorem}

\begin{proof}
Write $A=V_{\R}/\Lambda$, with Riemann form $E$, and choose $\tau$ such that $\mathcal O=\Z[\tau]$. Set $S=\mathcal O\Lambda=\Lambda+\tau\Lambda$ and $b=[S:\Lambda]$. The form $E'=fE$ is integral on $S$. Indeed the cross terms use the integral endomorphism $f\tau$, and Rosati compatibility gives
\[
 E(\tau x,\tau y)=\operatorname{Nm}(\tau)E(x,y).
\]
The quotient with lattice $S$ therefore has polarization volume
\[
 \Delta'=\frac{f^{2n}\Delta}{b}.
\]
On $S$ use the integral Hermitian form
\[
 H(x,y)=E'(\tau x,y)-\bar\tau E'(x,y),
 \qquad
 E'(x,y)=\operatorname{Tr}_{K/\Q}
          \frac{H(x,y)}{\tau-\bar\tau}.
\]
The first formula is linear in $x$ and conjugate-linear in $y$. The codifferent of $\mathcal O$ is $(\tau-\bar\tau)^{-1}\mathcal O$. Consequently a perfect integral Hermitian pairing gives a perfect alternating $\Z$-pairing by the second formula. This assertion holds for projective modules and requires no choice of a free basis.

Choose a maximal isotropic $K$-subspace $W$ of dimension $n$. The module $P_0=S\cap W$ is saturated: if a nonzero $a\in\mathcal O$ has $ax\in W$, then $x\in W$. The quotient $S/P_0$ is torsion-free over a Dedekind domain, hence projective. Choose a splitting
\[
 S=P_0\oplus Q.
\]
Nondegeneracy identifies $W$ with the conjugate dual of $Q_K$ by $w\mapsto H(w,\cdot)|_Q$. Let $P'$ be the inverse image of $Q^\vee$. Integrality gives $P_0\subset P'$. On $\Gamma=P'\oplus Q$, the form becomes
\[
 H(p+q,p'+q')=p(q')+\overline{p'(q)}+C(q,q'),
 \qquad C=H|_{Q\times Q}.
\]
The off-diagonal evaluation pairing is perfect, so this Hermitian form is unimodular. Its alternating form is therefore principal. Comparing the polarization volumes on $S$ and $\Gamma$ gives
\[
 [\Gamma:S]=\Delta'.
\]

Choose $\eta\in\mathcal O$ with $\eta+\bar\eta=t_K$, and let $c:Q\to Q^\vee$ be the $\mathcal O$-linear map $c(q)(q')=C(q,q')$. Enlarge the lattice once more, to
\[
 \Gamma_t=\frac1{t_K}P'\oplus Q,
\]
and give it the form $t_KH$. The map
\[
 Q^\vee\oplus Q\longrightarrow\Gamma_t,
 \qquad (p,q)\longmapsto\frac p{t_K}+q-
                                  \frac{\eta c(q)}{t_K}
\]
is an integral module isomorphism. Its off-diagonal pairing is evaluation; its pairing on the second summand is
$t_KC-\eta C-\bar\eta C=0$. It is thus an isometry from $H(Q)$.
In particular this is a principal hyperbolic lattice and
\[
 [\Gamma_t:S]=t_K^{2n}\Delta'.
\]
All these lattices lie in the same real vector space with the same complex structure. The Riemann form has only been multiplied by a positive integer, so positivity is preserved. The inclusion $\Lambda\subset\Gamma_t$ defines a holomorphic $K$-linear isogeny with
\[
 q^*\Theta=t_K fL,
 \qquad \deg q=b\,t_K^{2n}\Delta'=(t_K f)^{2n}\Delta.
\]

For an effective detecting $n$-cycle $Z$ on $P$, finite flat pullback and the projection formula give
\[
 \deg_L(q^*Z)=\frac{\deg q}{(t_K f)^n}\deg_\Theta Z
             =t_K^n f^n\Delta\deg_\Theta Z.
\]
The pullback has nonzero Weil projection because a $K$-linear isogeny identifies the rational Weil spaces. Some irreducible component is a detector, and its degree is at most the total degree.

Finally every projective rank-$n$ module over $\mathcal O$ has Steinitz form $\mathcal O^{n-1}\oplus I$. Isomorphic modules give isometric hyperbolic modules, so at most $h_K$ types are required. For each type the compatible period domain is connected; a neat arithmetic quotient gives its principal family. Passing to a finite level cover does not affect the bound.

Choose the representative $I$ integral, of norm $m$. In coordinates,
\[
 Q=\mathcal O^{n-1}\oplus I,
 \qquad Q^\vee=\mathcal O^{n-1}\oplus\bar I^{-1},
 \qquad m\bar I^{-1}=I.
\]
Thus $(p,q)\mapsto(mp,q)$ embeds $H(Q)$ into the standard lattice $\mathcal O^{2n}$ and multiplies its Hermitian form by $m$. Transporting the complex structure through this map gives an isogeny $P_I\to P_U$ with pullback polarization $m\Theta_I$. Since both polarizations are principal, its degree is $m^{2n}$. Pullback of a middle-dimensional detector therefore costs $m^n$. Taking the smallest norm in each ideal class proves the bound involving $m_K$. All isogenies here preserve CM.

For the numerical estimate, an integral ideal $J$ embedded in $\C$ has covolume $\operatorname{Nm}(J)\sqrt{|\operatorname{disc}K|}/2$. The planar convex-body theorem applied to a disk gives $0\ne\alpha\in J$ with
\[
 |\operatorname{Nm}(\alpha)|\le
 \frac2\pi\sqrt{|\operatorname{disc}K|}\operatorname{Nm}(J).
\]
Indeed a centrally symmetric convex set of area exceeding four times the covolume contains a nonzero lattice point: two points of its half-sized copy have the same image in a lattice fundamental domain, and their difference lies in the original set. Apply this to disks with radii decreasing to the stated bound. The ideal $(\alpha)J^{-1}$ is integral, lies in the inverse class of $J$, and has norm at most $2\sqrt{|\operatorname{disc}K|}/\pi$. This proves the estimate for every class.
\end{proof}

For a free maximal-order lattice with $t_K=2$, suppose in addition that its Hermitian form is even. The restriction $C$ in the proof has even diagonal. In a basis of $Q$ there is then an integral map $R:Q\to Q^\vee$ with $C(q,q')=R(q)(q')+\overline{R(q')(q)}$: choose the diagonal entries by the trace equation and choose one of each pair of off-diagonal entries to be zero. The shear $q\mapsto q-R(q)$ already makes the lattice $\Gamma$ hyperbolic, without the final scaling. The same proof gives the improved cost $f^n\Delta$ whenever the saturated maximal-order lattice with polarization $fE$ is free and even.

The theorem removes dependence on arbitrary integral polarization types and ideal classes from the existence question. It leaves one principal hyperbolic supplier for each field and dimension. It does not construct that supplier.
